% package
\usepackage{alltt}
\usepackage{amsfonts}
\usepackage{amsmath}
\usepackage{amssymb}
\usepackage{amsthm}
\usepackage{bm}
\usepackage{booktabs}
\usepackage{caption}
\usepackage{color}
\usepackage[colorlinks]{hyperref}
\usepackage{enumitem}
\usepackage{fancyhdr}
\usepackage{geometry}
\usepackage{graphicx}
\usepackage{mathdots}
\usepackage{mathtools}
\usepackage{microtype}
\usepackage{multirow}
\usepackage{multicol}
\usepackage{pdflscape}
\usepackage{pgffor}
\usepackage{pgfplots}
\usepackage{siunitx}
\usepackage{slashed}
\usepackage{tabularx}
\usepackage{tikz}
\usepackage{tkz-euclide}
\usepackage[normalem]{ulem}
\usepackage{xcolor}
\usepackage{xparse}
\usepackage{xpatch}
\usepackage[all]{xy}
\usepackage{imakeidx}

% layout
\pagestyle{fancy}
\fancyhf{}
\fancyfoot[C]{\thepage}
\lhead{\emph{\nouppercase{\leftmark}}}
\setlength{\headheight}{14.5pt}

% title
\author{ioshift}
\date{\today}
\title{\Huge{\bf{\ntitle}}}

% Others
\pgfplotsset{compat=1.18}
\setlength{\parindent}{0em}
\setlength{\tabcolsep}{0pt}

%%%%%%%%%%%%%%%%%%%%%%%%%%
%%%%%%%%%% math %%%%%%%%%%
%%%%%%%%%%%%%%%%%%%%%%%%%%

% Theorems
\theoremstyle{definition}
\newtheorem{thm}{Theorem}[section]
\newtheorem{axiom}[thm]{Axiom}
\newtheorem{conjecture}[thm]{Conjecture}
\newtheorem{cor}[thm]{Corollary}
\newtheorem{eg}[thm]{Example}
\newtheorem{lemma}[thm]{Lemma}
\newtheorem{prop}[thm]{Proposition}

\newtheorem*{question}{Question}
\newtheorem*{aim}{Aim}
\newtheorem*{assumption}{Assumption}
\newtheorem*{claim}{Claim}
\newtheorem*{ex}{Exercise}
\newtheorem*{fact}{Fact}
\newtheorem*{law}{Law}
\newtheorem*{remark}{Remark}
\newtheorem*{rrule}{Rule}
\newtheorem*{warning}{Warning}

\newenvironment{solution}{\begin{proof}[Solution]}{\end{proof}}

% Special sets
\newcommand{\C}{\mathbf{C}}
\newcommand{\R}{\mathbf{R}}
\newcommand{\Q}{\mathbf{Q}}
\newcommand{\Z}{\mathbf{Z}}
\newcommand{\N}{\mathbf{N}}
\newcommand{\F}{\mathbf{F}}

% Symbols
\newcommand{\noteq}{\not=}
\newcommand{\re}{\mathrm{Re\,}}
\newcommand{\im}{\mathrm{Im\,}}
\newcommand{\poly}[1]{\mathcal{P}(#1)}
\newcommand{\comp}[1]{#1^{\perp}}
\newcommand{\abs}[1]{\left\lvert #1\right\rvert}
\newcommand{\norm}[1]{\left\lVert #1\right\rVert}
\newcommand{\floor}[1]{\left\lfloor #1\right\rfloor}
\newcommand{\ceil}[1]{\left\lceil #1\right\rceil}
\newcommand{\set}[1]{\left\{#1\right\}}
\newcommand{\limit}[2]{\lim\limits_{#1\to #2}}
\newcommand{\bmap}[1]{\mathcal{B}(#1)} 
\newcommand{\lmap}[1]{\mathcal{L}(#1)}
\newcommand{\lmat}[1]{\mathcal{M}(#1)}
\newcommand{\iprod}[1]{\langle#1\rangle}
\newcommand{\by}[2]{$#1$-by-$#2$}
\newcommand{\trans}[1]{#1^\mathrm{t}}
\newcommand{\bil}[1]{#1^{(2)}}
\newcommand{\symbil}[1]{#1^{(2)}_{\mathrm{sym}}}
\newcommand{\altbil}[1]{#1^{(2)}_{\mathrm{alt}}}
\newcommand{\altmul}[2]{#1^{(#2)}_{\mathrm{alt}}}
\renewcommand{\span}{\operatorname{span}}
\renewcommand{\d}{\mathop{}\!\mathrm{d}} 
\renewcommand{\vec}[1]{\boldsymbol{\mathbf{#1}}}
\renewcommand{\emptyset}{\varnothing}
\DeclareMathOperator{\lcm}{lcm}
\DeclareMathOperator{\E}{E}
\DeclareMathOperator{\Var}{Var}
\DeclareMathOperator{\Null}{null\,}
\DeclareMathOperator{\range}{range\,}
\DeclareMathOperator{\volume}{volume\,}
\DeclareMathOperator{\tr}{tr}
\DeclareMathOperator{\perm}{perm}
\DeclareMathOperator{\sign}{sign}

%%% Automaic Definition Glossary
\newtheorem{definition}[thm]{Definition}
\newtheorem{notations}[thm]{Notation}
\newcounter{cnt}
\newcommand{\newterm}[3]{%
    \global\expandafter\unexpanded\expandafter{
    \expandafter\gdef\csname term#1\endcsname}{#2}%
    \global\expandafter\unexpanded\expandafter{
    \expandafter\gdef\csname cat#1\endcsname}{#3}%
}
\newenvironment{defi}[1]{
    \refstepcounter{cnt}\label{\thecnt}
    \newterm{\arabic{cnt}}{#1}{0}
    \begin{definition}\textbf{#1}
    
}{\end{definition}}
\newenvironment{notation}[1]{
    \refstepcounter{cnt}\label{\thecnt}
    \newterm{\arabic{cnt}}{#1}{0}
    \begin{notations}\textbf{#1}
    
}{\end{notations}}
\newcommand{\addelement}[1]{
    \refstepcounter{cnt}\label{\thecnt}%
    \newterm{\arabic{cnt}}{#1}{0}%
}
\newcommand{\addclass}[1]{
    \refstepcounter{cnt}\label{\thecnt}%
    \newterm{\arabic{cnt}}{#1}{1}%
}
% \let\oldsection\section
% \renewcommand{\section}[1]{
%   \oldsection{#1}\addclass{#1}
% }
% \titleformat{\section}{\normalfont\Large\bfseries}{}{0pt}{\thesection\quad}
\newcommand{\printglossary}{
    \newpage
    \section{Glossary}
    \foreach \x in {1,...,\arabic{cnt}} {
        \ifnum\csname cat\x\endcsname=0
            \hyperref[\x]{\csname term\x\endcsname}\par\addvspace{0.5em}
        \else
            \ifnum\x>1
                \vspace{1em}
            \fi
            {\large\textbf{\csname term\x\endcsname}\par\addvspace{0.7em}}
        \fi
    }
}